\documentclass[12pt, a4paper]{article}

% ── Packages ──────────────────────────────────────────────────────────────────
\usepackage[a4paper, margin=2.5cm]{geometry}
\usepackage{graphicx}
\usepackage{float}
\usepackage{amsmath}
\usepackage{amssymb}
\usepackage{xcolor}
\usepackage{mdframed}
\usepackage{parskip}
\usepackage{enumitem}
\usepackage{titlesec}
\usepackage{booktabs}
\usepackage{hyperref}
\usepackage[T1]{fontenc}
\usepackage{tikz}
\usepackage{listings}
\usetikzlibrary{arrows.meta}

% ── Graphics path ─────────────────────────────────────────────────────────────
\graphicspath{{../../Images_for_latex/Lecture_8/}}

% ── Colours ───────────────────────────────────────────────────────────────────
\definecolor{keyblue}{RGB}{30, 80, 160}
\definecolor{keybluebg}{RGB}{220, 235, 255}
\definecolor{noteorange}{RGB}{200, 100, 0}
\definecolor{noteorangebg}{RGB}{255, 240, 215}
\definecolor{warnred}{RGB}{180, 30, 30}
\definecolor{warnredbg}{RGB}{255, 220, 220}
\definecolor{defgreen}{RGB}{20, 110, 50}
\definecolor{defgreenbg}{RGB}{215, 245, 225}
\definecolor{headernavy}{RGB}{25, 35, 90}

% ── Box environments ──────────────────────────────────────────────────────────
\newmdenv[linecolor=keyblue,backgroundcolor=keybluebg,linewidth=1.5pt,
  leftline=true,rightline=false,topline=false,bottomline=false,
  innerleftmargin=10pt,innerrightmargin=8pt,
  innertopmargin=6pt,innerbottommargin=6pt,
  skipabove=8pt,skipbelow=8pt]{keybox}

\newmdenv[linecolor=noteorange,backgroundcolor=noteorangebg,linewidth=1.5pt,
  leftline=true,rightline=false,topline=false,bottomline=false,
  innerleftmargin=10pt,innerrightmargin=8pt,
  innertopmargin=6pt,innerbottommargin=6pt,
  skipabove=8pt,skipbelow=8pt]{notebox}

\newmdenv[linecolor=warnred,backgroundcolor=warnredbg,linewidth=1.5pt,
  leftline=true,rightline=false,topline=false,bottomline=false,
  innerleftmargin=10pt,innerrightmargin=8pt,
  innertopmargin=6pt,innerbottommargin=6pt,
  skipabove=8pt,skipbelow=8pt]{warnbox}

\newmdenv[linecolor=defgreen,backgroundcolor=defgreenbg,linewidth=1.5pt,
  leftline=true,rightline=false,topline=false,bottomline=false,
  innerleftmargin=10pt,innerrightmargin=8pt,
  innertopmargin=6pt,innerbottommargin=6pt,
  skipabove=8pt,skipbelow=8pt]{defbox}

% ── Section formatting ────────────────────────────────────────────────────────
\titleformat{\section}[block]{\large\bfseries\color{headernavy}}{}{0pt}{}
  [\vspace{2pt}\rule{\linewidth}{0.8pt}\vspace{4pt}]
\titleformat{\subsection}[block]{\normalsize\bfseries\color{keyblue}}{}{0pt}{}

% ── Document ──────────────────────────────────────────────────────────────────
\begin{document}

% ══════════════════════════════════════════════════════════════════════════════
\clearpage
\section*{Exercise: BLE RSSI Trilateration in Python}
% ══════════════════════════════════════════════════════════════════════════════
\vspace{-0.3cm}
\subsection*{BLE RSSI}
\begin{warnbox}
We specfically want to simulate noisy RSSI measurements in a 5\,$\times$\,10\,m room,
invert them to distances, and find the position by optimisation.

Parameters follow Bembenik \& Falcman (2020): $\gamma = 1.74$,
$f = 2.45$\,GHz, TxPower $A = -10\gamma\log_{10}(4\pi f/c)
\approx -35$\,dBm, $\sigma_S = 3$\,dBm RSSI noise.
\end{warnbox}

The idea behind BLE positioning is straightforward. Small BLE beacons are fixed around the room at positions that are known in advance — measured once at setup and stored.

They just continuously broadcast a radio signal; they do not know where the receiver is.
The receiver (e.g. a smartphone) listens to all beacons it can hear. 
Because radio signals weaken with distance — the further the beacon, the weaker the signal that arrives — 
the receiver uses the measured signal strength from each beacon as a proxy for distance

\subsection*{BLE RSSI Distance Model}
When a radio signal travels distance $\rho$ through space it loses power
according to the \textbf{free-space path loss} model. In free space the
signal spreads out in all directions ---  the further
it travels, the more the received
power drops. This is captured by:
\[
  \ell = \left(\frac{4\pi\rho}{\lambda}\right)^{\!\gamma}
       = \left(\frac{4\pi\rho f}{c}\right)^{\!\gamma}
\]
The exponent $\gamma$ controls how fast the signal decays with distance.
\begin{itemize}
  \item In free space $\gamma = 2$, meaning power falls as $1/\rho^2$.
  \item In practice $\gamma$ must be calibrated for the environment;
we use $\gamma = 1.74$ taken from Bembenik \& Falcman (2020), who
measured it experimentally in an office room.
\end{itemize}
\begin{warnbox}
Interestingly this is below the free-space value of 2 --- one might
expect walls to absorb signal and cause faster decay, but reflections
can constructively add energy toward the receiver, slowing it instead.
\end{warnbox}
Taking the path loss into the log domain (converting to dBm) gives the
RSSI formula, here added with noise:
\[
  \boxed{S = \underbrace{-10\gamma\log_{10}\rho}_{\text{decreases as }\rho\text{ grows}}
           + \underbrace{10\gamma\log_{10}\!\left(\frac{4\pi f}{c}\right)}_{A\;=\;\text{constant at 1\,m}} + \nu_i} \quad \text{Where} \quad
  \nu_i \sim \mathcal{N}(0,\sigma_S^2)
\]
The first term is what carries the distance information --- as $\rho$
increases, $S$ becomes more negative (weaker signal).

$A$ is a constant that depends only on frequency and $\gamma$; it is
calibrated at 1\,m by the manufacturer and stored in the beacon.

The noise $\nu_i \sim \mathcal{N}(0, \sigma_S^2)$ lumps together
two sources of error.
First, small electronic and quantisation noise from the \textbf{sensor}
itself.
More significantly, \textbf{multipath}: the signal bounces off walls
and furniture so multiple copies arrive simultaneously; depending on
their phase they add to or subtract from the direct signal, shifting
the RSSI reading even when the receiver is perfectly stationary.


\subsection*{Signal inversion}
Rearranging this, result in us being able to get distance from a measured RSSI reading:
\[
  \hat{\rho}_i = 10^{\,\frac{A - S_i}{10\gamma}}
\]

\begin{center}
\renewcommand{\arraystretch}{1.4}
\begin{tabular}{p{1.3cm} p{3.2cm} p{7.5cm}}
  \toprule
  \textbf{Symbol} & \textbf{Name} & \textbf{Meaning} \\
  \midrule
  $\ell$ & Path loss & Ratio of transmitted power to received power
    (dimensionless, $> 1$).  Larger $\ell$ means more signal lost. \\
  $S$       & RSSI [dBm]         & Received signal strength; more negative = weaker = further away \\
  $\rho$    & Range [m]          & Distance from beacon to receiver \\
  $\gamma$  & Path-loss exponent & $\approx 2$ free space; 1.6--1.8 indoors (paper measured 1.736) \\
  $f$       & Frequency [Hz]     & BLE: $\approx 2.4$\,GHz \\
  $\lambda$ & Wavelength & $\lambda = c/f$; for BLE at 2.4\,GHz,
    $\lambda \approx 12.5$\,cm. \\
  $A$       & TxPower [dBm]      & RSSI at 1\,m; calibrated by manufacturer and stored in beacon \\
  $c$       & $3\times10^8$\,m/s. & Speed of light, electromagnetic waves travel at this speed \\
  \bottomrule
\end{tabular}
\end{center}
\subsection*{Why not just time-of-flight?}

\begin{notebox}
  \textbf{Bluetooth --- TOF impossible on standard hardware.}
  Radio travels at $c = 3\times10^8$\,m/s, so a 10\,m gap takes only
  33\,ns to cross. Standard BLE chips have processing delays of
  $\sim$150\,\textmu s before software even registers a packet arrived
  --- that delay alone gives a fictitious range of:
  \[
    \frac{150\times10^{-6}\,\text{s} \times 3\times10^8\,\text{m/s}}{2}
    = 22{,}500\,\text{m}
  \]
  The travel time is completely buried in the processing overhead, so
  RSSI is used instead --- it costs nothing extra since every BLE chip
  measures received power as part of normal operation.
\end{notebox}

\begin{notebox}
  \textbf{Games-on-Track and GPS --- TOF works because of dedicated hardware
  and is more accurate.}
  Systems like Games-on-Track and GPS use chips built specifically with
  nanosecond-precision timestamping circuits, making TOF both possible
  and far more accurate than RSSI.

  Rather than measuring signal strength, both the transmitter and receiver
  generate the \textbf{same known pattern} --- a chirp (Games-on-Track)
  or a pseudo-random code (GPS).

  The receiver then computes the \textbf{cross-correlation} between its
  local copy and the incoming signal; the correlation peaks sharply at
  exactly the time shift where the two align, giving the travel time
  $\Delta t$ and thus the range $r = c\,\Delta t$.

  This is far more robust indoors than RSSI because it locks onto the
  \textbf{first-arriving} copy of the signal --- the direct path ---
  and ignores later reflections off walls entirely.
\end{notebox}

\subsection*{Optimisation-based trilateration}
Because RSSI noise means the four estimated distances rarely place
the receiver at exactly one consistent point, we instead find the
position $(x, y)$ that is most consistent with all four measurements
simultaneously by minimising the sum of squared residuals:
\[
  \min_{x,y}\;\sum_{i=1}^{4}
  \Bigl[\underbrace{\sqrt{(x-b_{x,i})^2+(y-b_{y,i})^2}}_{\text{geometric dist to beacon }i}
  \;-\;\underbrace{\hat{\rho}_i}_{\text{estimated dist}}\Bigr]^2
\]
For each beacon $i$ the bracket computes the geometric distance from
the candidate position $(x,y)$ to the known beacon location
$(b_{x,i}, b_{y,i})$, minus the distance estimated from the RSSI
reading $\hat{\rho}_i$.

If $(x,y)$ were placed exactly on beacon $i$'s circle the residual
for that beacon would be zero --- but that point very likely lie far from
the other three circles, making their residuals large.

Summing over all four beacons prevents this: the solver is forced to
find the position that is simultaneously close to all circles at once,
which corresponds to the true receiver location.
This is solved numerically with \texttt{scipy.optimize.minimize}
(Nelder-Mead).

\subsection*{Results}

With $\sigma_S = 3$\,dBm and 40 trials per location the simulation gives:

\begin{center}
\renewcommand{\arraystretch}{1.4}
\begin{tabular}{lcc}
  \toprule
  \textbf{True position} & \textbf{Mean error} & \textbf{90th pct} \\
  \midrule
  (1.25, 2.50)\,m & 2.95\,m & 5.07\,m \\
  (2.50, 5.00)\,m & 2.27\,m & 4.46\,m \\
  (3.75, 7.50)\,m & 2.46\,m & 4.39\,m \\
  \bottomrule
\end{tabular}
\end{center}

Bembenik \& Falcman report a raw baseline (no filtering) of 4.62\,m
mean and 6.64\,m at the 90th percentile across two large irregular
office rooms spanning roughly 20--30\,m with walls between them.
Our errors are lower since the 5$\times$10\,m room with four ideal
corner beacons is a considerably simpler geometry --- no irregular
walls, no signal bleed from adjacent rooms, and beacons optimally
placed at corners.
\clearpage
\begin{notebox}
  \textbf{What the red dot clouds tell us.}
  Each red dot is one complete trilateration result --- the same
  position estimated from one noisy set of four RSSI readings.
  The spread of the cloud is the \emph{positioning noise floor}: even
  standing perfectly still, each trial lands somewhere different because
  RSSI fluctuates due to multipath and sensor noise.

  Unlike LiDAR where noise enters directly as small errors in range
  ($\sim$5\,cm) and angle ($\sim$1$^\circ$), here the noise is on the
  signal strength itself.
  As the beacon gets further away the signal weakens, so the same
  3\,dBm noise represents an increasingly large fraction of the total
  reading --- the signal-to-noise ratio degrades with distance.
  The log scale makes this worse: the RSSI curve flattens at large
  distances, meaning a small change in $S$ maps to a large jump in
  the estimated $\rho = 10^{(A-S)/(10\gamma)}$, translating 3\,dBm
  of noise into \textbf{metres} of range uncertainty.
\end{notebox}

\begin{warnbox}
  \textbf{Why the estimates scatter so much.}
  The signal from each beacon does not arrive via just one path.
  It bounces off walls, floor, and furniture, so \textbf{multiple
  copies} arrive simultaneously.
  RSSI measures the \emph{total} received power --- the direct copy
  plus all reflections combined.
  This blended power fluctuates moment to moment (a door opens,
  someone walks past) as reflections shift in and out of phase.
  Fluctuating RSSI gives a fluctuating range estimate, hence a
  different trilateration result each time.
  Typical indoor accuracy: \textbf{1--5\,m}.
\end{warnbox}

\end{document}